\clearpage
\phantomsection
\addcontentsline{toc}{chapter}{KATA PENGANTAR}
\begin{center}
    \textbf{\large KATA PENGANTAR}\\[3em]
\end{center}
%-----------------------------------------

Puji syukur kehadirat Allah SWT atas berkat rahmat dan karunia-Nya, penulis dapat melaksanakan dan menyelesaikan {\tipe} di {\tempat}. Laporan ini disusun dalam rangka memenuhi sebagian persyaratan akademik pada {\prodi}, {\departemen}, {\fakultas}, {\universitas}.

{\tipe} dengan judul ``{\judulid}'' ini disusun berdasarkan kegiatan kerja praktik, observasi lapangan, studi literatur, serta analisis teknis yang berkaitan dengan sistem PLTS terapung. Secara khusus, laporan ini membahas kelayakan dan pola degradasi modul PV pasca-terendam di PLTS Terapung Cirata menggunakan pengujian kurva I--V dan indikator arus--tegangan.

Pelaksanaan kerja praktik ini memberikan kesempatan kepada penulis untuk memahami kegiatan operasi dan pemeliharaan pada PLTS terapung, terutama yang berkaitan dengan kondisi abnormal pada modul PV. Melalui laporan ini, penulis berupaya menghubungkan fenomena modul PV pasca-terendam dengan parameter elektrik hasil pengujian, seperti $P_{max}$, $V_{oc}$, $I_{sc}$, $V_{mp}$, $I_{mp}$, Fill Factor, serta indikator $R_{VM}$ dan $R_{IM}$. Analisis tersebut digunakan sebagai dasar awal untuk membaca perubahan performa modul, mengidentifikasi kecenderungan degradasi, serta menyusun arah pemeriksaan teknis yang lebih terstruktur pada kegiatan \textit{operation and maintenance}.

{\tipe} ini dapat diselesaikan tidak lepas dari bantuan, arahan, bimbingan, dan kerja sama dari berbagai pihak. Berkenaan dengan hal tersebut, penulis menyampaikan ucapan terima kasih kepada yang terhormat:

\begin{enumerate}
    \item Bapak {\koordepartemen}, selaku Ketua {\departemen}, {\fakultas}, {\universitas}.
    
    \item Bapak {\kaprodi}, selaku Ketua {\prodi}, {\departemen}, {\fakultas}, {\universitas}.
    
    \item Bapak {\pembimbing}, selaku Dosen Pembimbing Kerja Praktik, yang telah memberikan arahan, bimbingan, dan masukan selama pelaksanaan kerja praktik serta penyusunan laporan ini.
    
    \item Bapak {\pembimbingPerusahaan}, selaku {\jabatanPembimbing} di {\tempat}, yang telah memberikan arahan, informasi, dan kesempatan kepada penulis untuk memahami kegiatan operasi dan pemeliharaan di lingkungan PLTS Terapung Cirata.
    
    \item Segenap staf dan karyawan {\tempat} yang telah membantu penulis selama pelaksanaan kerja praktik.
    
    \item Orang tua dan keluarga penulis yang selalu memberikan doa, dukungan, dan motivasi selama pelaksanaan kerja praktik serta penyusunan laporan ini.
    
    \item Seluruh pihak yang tidak dapat penulis sebutkan satu per satu, yang telah memberikan bantuan dan dukungan dalam pelaksanaan kerja praktik dan penyusunan laporan ini.
\end{enumerate}

Penulis menyadari bahwa laporan ini masih memiliki kekurangan. Oleh karena itu, penulis mengharapkan kritik dan saran yang membangun agar laporan ini dapat menjadi lebih baik. Akhirnya, semoga segala bantuan yang telah diberikan oleh semua pihak menjadi amal kebaikan dan semoga laporan kerja praktik ini dapat memberikan manfaat bagi pembaca, pihak akademik, maupun pihak lain yang membutuhkan kajian teknis mengenai evaluasi performa modul PV pada PLTS terapung.

\begin{flushright}
    Yogyakarta, \tglpengesahan\\[1.25cm]
    \begin{tabular}{c@{\hskip 1cm}c}
        \underline{\penulisSatu} & \underline{\penulisDua} \\
        \nimSatu & \nimDua
    \end{tabular}
\end{flushright}